\chapter{Dataset Preparation}
This chapter details the data pipeline designed to generate the expert trajectories used to train and evaluate the Flow Planner framework. The process begins with the selection of a suitable evaluation environment, followed by the training of an expert agent using Proximal Policy Optimization (PPO) to generate the demonstration dataset. The datasets as well as the trained agents can be found on my personal \href {https://huggingface.co/frankcholula}{Hugging Face Hub}.

\section{Environment Suite}
To test the capabilities of a generative planner, a suite of three distinct environments from the Gymnasium \cite{towers2024gymnasium} Library was selected, each representing a unique set of challenges. While this dissertation will focus on the empirical evaluation of LunarLanderContinuous-v3 for a more focused analysis, the other environments are included here for future research into scaling the Flow Planner framework to more complex locomotion and vision-based tasks.
\subsection{LunarLanderContinuous-v3}
This environment presents a classic control problem where the agent's objective is to land a craft safely and with minimal fuel consumption between two designated flags. The low-dimensional state and continuous dynamics makes it an ideal benchmark for testing an agent’s ability to learn precise thrust control and solve a delayed credit assignment problem, where early actions have critical consequences. A successful landing requires scoring over 200 points, which incorporates rewards for proximity, velocity, stability, and fuel efficiency, with a bonus of +100 for a safe landing and -100 for crashing.
\begin{figure}[H]
  \centering
  \includegraphics[width=1.0\textwidth]{lunar_lander}
  \caption{LunarLanderContinuous-v3 Environment}
  \label{fig:lunar}
\end{figure}


\subsection{BipedalWalker-v3}
This environment increases the complexity by focusing on continuous locomotion over a randomly generated, uneven terrain. Success requires the agent to master the coordination of multiple joints to develop a robust and efficient walking gait, making it a challenging benchmark for learning stable locomotion policies.
\begin{figure}[H]
  \centering
  \includegraphics[width=1.0\textwidth]{bipedal_walker}
  \caption{BipedalWalker-v3 Environment}
  \label{fig:bipedal}
\end{figure}


\subsection{CarRacing-v3}
This environment introduces the challenge of vision-based control from high-dimensional 96x96 pixel images. The objective is to complete a lap on a procedurally generated track, testing the agent’s representation learning capabilities and its ability to perform high-speed continuous control from raw pixels. The detailed procedure for processing these visual observations and generating the final latent dataset is provided in Appendix \ref{app:latent_dataset}.
\begin{figure}[H]
  \centering
  \includegraphics[width=1.0\textwidth]{car_racing}
  \caption{CarRacing-v3 Environment}
  \label{fig:carracing}
\end{figure}

\section{Expert Policy Training}
\label{sec:expert_training}
To generate the expert datasets, a Proximal Policy Optimization (PPO) agent was trained in each environment until convergence. PPO was chosen for its exceptional training stability. As a state-of-the-art on-policy algorithm, it optimizes the policy by taking a series of small, conservative steps. It uses a “clipped” surrogate objective that acts as a guardrail, preventing any single update from changing the policy too dramatically \cite{schulmanProximalPolicyOptimization2017}. This inherent stability makes PPO a robust workhorse for the continuous control challenges in these environments.

All agents were trained using the tuned hyper-parameters supplied by the RL Baselines3 Zoo \cite{rl-zoo3} library, with the final performance summarized in Table \ref{tab:ppo_llc}, \ref{tab:ppo_bipedal} and \ref{tab:ppo_carracing}.
% --- LunarLanderContinuous-v3 ---
\begin{table}[H]
\centering
\caption{PPO expert (LunarLanderContinuous-v3).}
\label{tab:ppo_llc}
\begin{tabular}{|l|p{7.5cm}|}
\hline
\textbf{Mean Reward}       & 236.78 $\pm$ 79.67 \\ \hline
\textbf{Wrappers}          & None \\ \hline
\textbf{Observation Space} & Box(8,) float32 \\ \hline
\textbf{Action Space}      & Box(2,) float32 \\ \hline
\end{tabular}
\end{table}

% --- BipedalWalker-v3 ---
\begin{table}[H]
\centering
\caption{PPO expert (BipedalWalker-v3).}
\label{tab:ppo_bipedal}
\begin{tabular}{|l|p{7.5cm}|}
\hline
\textbf{Mean Reward}       & 271.87 $\pm$ 2.08 \\ \hline
\textbf{Wrappers}          & None \\ \hline
\textbf{Observation Space} & Box(24,) float32 \\ \hline
\textbf{Action Space}      & Box(4,) float32 \\ \hline
\end{tabular}
\end{table}

% --- CarRacing-v3 ---
\begin{table}[H]
\centering
\caption{PPO expert (CarRacing-v3).}
\label{tab:ppo_carracing}
\begin{tabular}{|l|p{7.5cm}|}
\hline
\textbf{Mean Reward}       & 248.24 $\pm$ 168.67 \\ \hline
\textbf{Wrappers}          & \textbf{FrameSkip, Resize, GrayScale, FrameStack} \\ \hline
\textbf{Observation Space} & Box(64, 64, 2) uint8 \\ \hline
\textbf{Action Space}      & Box(3,) float32 \\ \hline
\end{tabular}
\end{table}

To ensure a comprehensive set of expert policies, agents were also trained using the A2C and TD3 algorithms. Although these methods demonstrated superior performance in certain environments, solely the PPO agents were used to generate the expert datasets. This decision was made to prioritize the quality and stability of the expert behavior for imitation. PPO's on-policy nature and clipped objective function consistently produce smooth and generalizable policies, which serve as a more robust and reliable foundation for training a generative model. While off-policy methods like TD3 can be highly performant, they can also yield more exploitative or brittle policies that are less suitable as a behavioral prior. 
For the benefit of future research and to ensure reproducibility, all trained models, including the A2C and TD3 agents, have been made publicly available on my \href {https://huggingface.co/frankcholula}{Hugging Face Hub}.
\section{Expert Dataset Collection}
To train the generative planners described in Chapter 3, a dataset of expert trajectories were required for each environment. This data was collected by executing the trained policies from Section \ref{sec:expert_training}. The Minari \cite{minari} Library’s  DataCollector wrapper was used to systematically record trajectory data during these rollouts. Approximately 300,000 to 350,000 steps were sampled for each environment to create a robust dataset capturing a wide range of expert behaviors. The resulting datasets are summarized in the table below. Note the number of episodes collected for each environment was tailored to its specific characteristics. For instance, the stochastic nature of LunarLander-v3 necessitated a large number of episodes (1000) to ensure a diverse set of successful landing trajectories. On the other hand, the high-dimensional visual observations in CarRacing-v3 resulted in a substantial dataset size from fewer but longer episodes (250). The dataset statistics for different environments are summarized in Table \ref{tab:datasets}.
\begin{table}[H]
\centering
\caption{Dataset statistics for different environments.}
\label{tab:datasets}
\begin{tabular}{|l|c|c|c|}
\hline
\textbf{Environment} & \textbf{LunarLanderContinuous-v3} & \textbf{BipedalWalker-v3} & \textbf{CarRacing-v3} \\
\hline
\textbf{\# Episodes} & \textbf{1000} & 500 & 250 \\
\hline
\textbf{\# Steps}    & 335K & \textbf{346K} & 245K \\
\hline
\textbf{Size}        & 33.7 MB & 49.9 MB & \textbf{446.8 MB} \\
\hline
\end{tabular}
\end{table}

Again, to ensure full reproducibility and to facilitate future research, these datasets have been version-controlled and are publicly available on my personal \href {https://huggingface.co/frankcholula}{Hugging Face Hub}.

\chapter{Experimental Results}
\label{sec:exp_results}
This chapter presents the empirical evaluation of the proposed Flow Planner Framework, with experiments designed to validate the performance of the hierarchical planning methodology and to analyze the impact of different architectural choices and conditioning signals. To ensure fair and reproducible comparisons, all experiments were conducted on a workstation equipped with an NVIDIA RTX 4080 Super GPU and an AMD Ryzen 7 7800X3D CPU. The LunarLander-v3 environment serves as the primary testbed for this analysis.

This chapter is structured as follows: Section \ref{sec:architecture} details and compares the performance of different neural network architectures for both flow matching and diffusion-based planners[cite diffuser here]. Leveraging the best-performing architecture from this initial analysis, Section \ref{sec:joint_planner_baseline}	 establishes a performance baseline by evaluating the monolithic Joint Planner under two different conditioning schemes. Finally, Section \ref{sec:flat_hieararchical_planning} evaluates the core contribution of this work: the decoupled, hierarchical planner that combines the Waypoint and Action Planners. Each of these experimental sections will present both qualitative analysis of  the open-loop trajectory generation and a qualitative, closed-loop MPC evaluation to measure the final planner’s performance via cumulative rewards over 10 episodes.

Unless otherwise specified, all models for the remainder of Section \ref{sec:exp_results} were trained for 100 epochs with a learning rate of 1e-3. For CNN-based models, the kernel size was standardized to 5. For the diffusion model comparison, models were trained with 1000 diffusion steps and sampled using 100 inference steps with a standard DDPM scheduler.

\section{Neural Network Architecture Comparisons}
The initial phase of this empirical study is to determine the most effective architectural backbone. A generative planner’s performance is directly tied to its network architecture, especially for sequential data such as a trajectory that possesses a strong temporal structure. This section establishes a simple MLP as a baseline, then evaluates its performance against two more sophisticated architectures: a 1D CNN \cite{osheaIntroductionConvolutionalNeural2015} and a U-Net \cite{ronnebergerUNetConvolutionalNetworks2015a}. Each architecture is implemented for both our novel flow matching planners and the existing diffusion-based models (used as baselines), to determine the optimal neural network backbone for the planning experiments that follow. Note that all models in this section are trained on the unconditional tasks of modeling the full state-action trajectories.

\label{sec:architecture}
\subsection{MLP Baseline}
As a foundational baseline, a simple Multi-Layer Perceptron (MLP) is used. The network first processes a flattened representation of the entire trajectory and projects it into a hidden dimension. This is then combined with a sinusoidal time embedding before being passed through a series of feedforward layers \cite{vaswaniAttentionAllYou2023}. While computationally efficient, this architecture deliberately ignores the inherent temporal coherence of the data. Nevertheless, it serves as a crucial baseline to quantify the performance gains achieved by more structured, sequence-aware architectures. The detailed implementation is provided below.

The MLP baseline is a fully-connected neural network designed to process the entire trajectory as a single, flattened vector. The architecture consists of three main components: an initial projection, a time embedding module, and a main processing block.
\begin{enumerate}
\item
\end{enumerate}


\subsection{1D CNN}
\subsection{U-Net}
\subsection{Results and Analysis}

\section{Planner Baseline: The Joint Planner}
\label{sec:joint_planner_baseline}
\subsection{Experimental Setup}
\subsection{Conditioning Signals and Fusion Strategy}
\subsection{Classifier-free Guidance vs. Modeling Conditional Distribution}
\subsection{Results and Analysis}

\section{Flat vs. Hierarchical Planning}
\label{sec:flat_hieararchical_planning}
\subsection{Individual Planner Performance}
\subsection{Hierarchical Plannning}


